\section{Kategorie dostępne i lokalnie prezentowalne}\label{cat:sec:accessible-presentable}

Konstrukcje uniwersalne można badać przez ich zachowanie względem małych fragmentów obiektów. Na przykład każdy zbiór jest kogranicą filtrowaną swoich skończonych podzbiorów, a każdy moduł jest kogranicą filtrowaną modułów zadanych przez skończenie wiele generatorów i relacji. W drugim przykładzie morfizmy diagramu nie muszą być monomorfizmami. Teoria kategorii dostępnych ujmuje obie sytuacje za pomocą zachowywania odpowiednich kogranic przez funktory reprezentowalne.

W tym rozdziale pracujemy z kategoriami zwykłymi, lokalnie małymi względem ustalonego uniwersum. Wszystkie rozważane diagramy indeksujące kogranice są małe w tym uniwersum. Korzystamy z aksjomatu wyboru i utożsamiamy liczby kardynalne z odpowiadającymi im początkowymi liczbami porządkowymi.

\subsection{Regularne liczby kardynalne i diagramy filtrowane}

\begin{definicja}[Regularność]\label{cat:def:accessible-regular}
Nieskończoną liczbę kardynalną $\kappa$ nazywamy regularną, gdy jej kofinalność jest równa $\kappa$. Równoważnie: dla każdej rodziny $(S_i)_{i\in I}$ spełniającej $|I|<\kappa$ i $|S_i|<\kappa$ dla wszystkich $i$ zachodzi
\[
 \left|\bigcup_{i\in I}S_i\right|<\kappa.
\]
\end{definicja}
Aby sprawdzić równoważność, przy regularnej $\kappa$ liczby $|S_i|$ mają wspólne ograniczenie kardynalne poniżej $\kappa$, gdyż ich supremum nie może być kofinalne w $\kappa$. Jeśli $\kappa=\aleph_0$, chodzi o sumę skończenie wielu skończonych zbiorów. Dla nieprzeliczalnej $\kappa$ oszacowanie przez iloczyn $\max(\aleph_0,|I|,\sup_i|S_i|)$ z samym sobą daje liczbę mniejszą od $\kappa$. Odwrotnie, kofinalna rodzina mniej niż $\kappa$ liczb porządkowych mniejszych od $\kappa$ dawałaby przedstawienie $\kappa$ jako sumy zbiorów o mocach mniejszych od $\kappa$.

W dalszym ciągu $\kappa$ jest zawsze nieskończoną regularną liczbą kardynalną. W szczególności $\aleph_0$ jest regularna.

\begin{definicja}[$\kappa$-filtrowanie]\label{cat:def:accessible-filtered}
Kategoria $I$ jest $\kappa$-filtrowana, gdy każdy diagram $D:J\to I$, w którym $|\operatorname{Mor}J|<\kappa$, ma kostożek w $I$. Uwzględniamy diagram pusty, zatem $I$ musi być niepusta.
\end{definicja}

Dla $\kappa=\aleph_0$ otrzymujemy zwykłe filtrowanie z definicji~\ref{cat:def:filtered}. W przypadku posetu, traktowanego jako kategoria, warunek oznacza dokładnie istnienie wspólnego ograniczenia górnego każdej rodziny mniej niż $\kappa$ elementów. Dla kategorii ogólnej potrzebna jest ponadto możliwość uzgadniania równoległych morfizmów.

\begin{stwierdzenie}[Kryterium $\kappa$-filtrowania]\label{cat:prop:accessible-filtered-criterion}
Kategoria $I$ jest $\kappa$-filtrowana wtedy i tylko wtedy, gdy spełnia następujące dwa warunki:
\begin{enumerate}
 \item każda rodzina mniej niż $\kappa$ obiektów (również pusta) ma wspólny obiekt docelowy, tj. istnieją morfizmy ze wszystkich jej obiektów do jednego obiektu;
 \item dla każdej rodziny mniej niż $\kappa$ par równoległych morfizmów $f_\alpha,g_\alpha:A_\alpha\to B$ o wspólnym celu istnieje $u:B\to C$ takie, że $uf_\alpha=ug_\alpha$ dla wszystkich $\alpha$.
\end{enumerate}
\end{stwierdzenie}
\begin{proof}
Pierwszy warunek wynika z zastosowania definicji do kategorii dyskretnej. Dla drugiego tworzymy kategorię indeksującą mającą jeden obiekt $b$ i po jednym obiekcie $a_\alpha$ dla każdej pary. Poza identycznościami ma dokładnie dwie strzałki $a_\alpha\to b$ dla każdego $\alpha$. Jej liczba morfizmów jest mniejsza od $\kappa$. Kostożek nad diagramem wysyłającym te strzałki na $f_\alpha,g_\alpha$ daje morfizm $u$ o żądanej własności.

Odwrotnie, niech $D:J\to I$ będzie diagramem o mniej niż $\kappa$ morfizmach. Liczba obiektów $J$ również jest mniejsza od $\kappa$, ponieważ każdy obiekt wyznacza inną identyczność. Z pierwszego warunku wybieramy obiekt $B$ i morfizmy $v_j:D(j)\to B$. Dla każdego $\alpha:j\to k$ w $J$ rozważamy parę
\[
 v_j,\quad v_kD(\alpha):D(j)\longrightarrow B.
\]
Liczba tych par jest mniejsza od $\kappa$. Drugi warunek daje $u:B\to C$ zerujący ich różnice w sensie równości złożeń. Rodzina $uv_j:D(j)\to C$ spełnia zatem warunek kostożka dla każdej strzałki $\alpha$.
\end{proof}

\begin{przyklad}[Poset małych podzbiorów]\label{cat:ex:accessible-small-subsets}
Dla zbioru $S$ niech
\[
 [S]^{<\kappa}=\{A\subseteq S:|A|<\kappa\}
\]
będzie posetem uporządkowanym przez zawieranie. Jest on $\kappa$-filtrowany: suma mniej niż $\kappa$ podzbiorów o mocy mniejszej od $\kappa$ znów ma moc mniejszą od $\kappa$ przez regularność. Kogranicą diagramu inkluzji $A\hookrightarrow S$ w $\Set$ jest $S$. Istotnie, każdy element należy do pewnego podzbioru jednoelementowego, a zgodność rodziny funkcji na małych podzbiorach oznacza, że jej wartość w danym elemencie nie zależy od wyboru podzbioru (można przejść do sumy dwóch podzbiorów).
\end{przyklad}

\subsection{Kogranice filtrowane w kategorii zbiorów}

Rozwijamy twierdzenie~\ref{cat:thm:filtered-finite} z rozdziału~\ref{sec:granice} na granice o mniej niż $\kappa$ morfizmach.

\begin{twierdzenie}[$\kappa$-filtrowane kogranice a małe granice]\label{cat:thm:accessible-filtered-limits}
Niech $I$ będzie małą $\kappa$-filtrowaną kategorią, $J$ małą kategorią o $|\operatorname{Mor}J|<\kappa$, a $D:I\times J\to\Set$ bifunktorem. Kanoniczna funkcja
\[
 \operatorname{colim}_{i\in I}\lim_{j\in J}D(i,j)
 \longrightarrow
 \lim_{j\in J}\operatorname{colim}_{i\in I}D(i,j)
\]
jest bijekcją.
\end{twierdzenie}
\begin{proof}
Dla każdego $i$ granicę po $J$ opisujemy jako zbiór rodzin $(x_j)_j$ z $x_j\in D(i,j)$ spełniających $D(i,\alpha)x_j=x_k$ dla $\alpha:j\to k$. Kanoniczna funkcja wysyła klasę takiej rodziny na rodzinę klas $([x_j])_j$. Jest dobrze określona, ponieważ morfizmy w $I$ działają współrzędnie i komutują z morfizmami w $J$.

Najpierw dowodzimy surjektywności. Niech $(\xi_j)_j$ będzie zgodną rodziną klas w kogranicach. Dla każdego $j$ wybieramy reprezentanta $x_j\in D(i_j,j)$. Ponieważ liczba obiektów $J$ jest mniejsza od $\kappa$, wybieramy wspólny obiekt $i$ z morfizmami $i_j\to i$ i przenosimy wszystkie reprezentanty do $D(i,j)$, zachowując ich oznaczenia. Dla każdej strzałki $\alpha:j\to k$ elementy $D(i,\alpha)x_j$ i $x_k$ mają tę samą klasę w kogranicy. Lemat~\ref{cat:lem:filtered-equality} daje morfizm $u_\alpha:i\to i_\alpha$ wyrównujący te dwa elementy. Jeżeli lemat początkowo daje dwie strzałki z $i$, stosujemy zwykłe filtrowanie, aby je uzgodnić i otrzymać jedną.

Wybieramy wspólny obiekt $t$ dla rodziny $(i_\alpha)_\alpha$ oraz $i$, z morfizmami $v_\alpha:i_\alpha\to t$ i $v:i\to t$. Przez stwierdzenie~\ref{cat:prop:accessible-filtered-criterion} istnieje $w:t\to t'$ wyrównujące jednocześnie wszystkie pary $v_\alpha u_\alpha,v:i\to t$. Morfizm $wv:i\to t'$ przenosi więc reprezentanty na rodzinę spełniającą wszystkie równania diagramu $J$. Rodzina ta reprezentuje pierwotny element prawej strony. Gdy $J$ jest pusta, wystarcza dowolny obiekt niepustej kategorii $I$ i jedyna pusta rodzina.

Dla injektywności weźmy dwie zgodne rodziny, odpowiednio w stopniach $i,i'$, o równych obrazach po prawej stronie. Przenosimy je do wspólnego stopnia $t$, otrzymując $(x_j)_j,(y_j)_j$. Dla każdego $j$ lemat~\ref{cat:lem:filtered-equality} daje morfizm z $t$ wyrównujący $x_j,y_j$. Jak w części o surjektywności, za pomocą wspólnego celu i jednoczesnego wyrównania mniej niż $\kappa$ par uzyskujemy jeden morfizm $t\to t'$ wyrównujący wszystkie współrzędne. Przeniesione rodziny są więc równe jako elementy granicy w stopniu $t'$. Zatem ich klasy po lewej stronie są równe. Dla pustego $J$ obie strony są jednoelementowe, ponieważ $I$ jest niepusta i dowolne dwa jej obiekty mają wspólny cel.
\end{proof}

\begin{przyklad}[Rola ograniczenia kardynalnego]\label{cat:ex:accessible-filtered-limit-failure}
Niech $I=[\kappa]^{<\kappa}$ i $J$ będzie dyskretną kategorią o zbiorze obiektów $\kappa$. Dla $A\in I$ oraz $j\in\kappa$ połóżmy $D(A,j)=\{0,1\}$, gdy $j\in A$, i $D(A,j)=\{0\}$ w przeciwnym razie. Morfizmy w $I$ działają przez inkluzje. Kogranica produktów identyfikuje się ze zbiorem funkcji $\kappa\to\{0,1\}$, których nośnik ma moc mniejszą od $\kappa$: każda funkcja w stopniu $A$ zeruje się poza $A$, a każda funkcja o małym nośniku należy do stopnia danego przez ten nośnik. Produkt kogranic jest natomiast całym zbiorem $\{0,1\}^\kappa$. Funkcja stale równa $1$ nie należy do obrazu porównania. Granice indeksowane przez dokładnie $\kappa$ obiektów nie muszą zatem komutować z $\kappa$-filtrowanymi kogranicami.
\end{przyklad}

\subsection{Obiekty prezentowalne}

\begin{definicja}[Obiekt $\kappa$-prezentowalny]\label{cat:def:accessible-presentable-object}
Niech $\mathcal C$ będzie lokalnie małą kategorią. Obiekt $P$ nazywamy $\kappa$-prezentowalnym, gdy funktor $\Hom_{\mathcal C}(P,-):\mathcal C\to\Set$ zachowuje wszystkie istniejące małe $\kappa$-filtrowane kogranice. Jeżeli $\mathcal C$ ma wszystkie takie kogranice, warunek oznacza, że dla każdego diagramu $D:I\to\mathcal C$ o $\kappa$-filtrowanej kategorii indeksującej funkcja
\begin{equation}\label{cat:eq:accessible-hom-colim}
 \operatorname{colim}_{i\in I}\Hom_{\mathcal C}(P,D(i))
 \longrightarrow \Hom_{\mathcal C}(P,\operatorname{colim}_I D)
\end{equation}
jest bijekcją. Obiekt $\aleph_0$-prezentowalny nazywamy skończenie prezentowalnym.
\end{definicja}

Zauważmy, że obiekt skończenie generowany i skończenie prezentowalny to na ogół różne pojęcia: samo generowanie nie kontroluje relacji pomiędzy generatorami.

\begin{stwierdzenie}[Kryterium faktoryzacji i równości]\label{cat:prop:accessible-presentable-factor}
Dla obiektu $P$ w kategorii mającej $\kappa$-filtrowane kogranice warunek prezentowalności jest równoważny następującym dwóm własnościom dla każdego takiego diagramu $D$ o kogranicy $L$ i morfizmach strukturalnych $\iota_i$:
\begin{enumerate}
 \item każdy $f:P\to L$ ma postać $f=\iota_i f_i$ dla pewnego $i$ i $f_i:P\to D(i)$;
 \item jeżeli $\iota_i f_i=\iota_j f_j$, to istnieją $u:i\to k$ i $v:j\to k$ takie, że $D(u)f_i=D(v)f_j$.
\end{enumerate}
\end{stwierdzenie}
\begin{proof}
Pierwsza własność jest dokładnie surjektywnością funkcji~\eqref{cat:eq:accessible-hom-colim}, ponieważ elementy kogranicy zbiorów są reprezentowane przez elementy poszczególnych składników. Druga własność jest dokładnie jej injektywnością przez lemat~\ref{cat:lem:filtered-equality}, zastosowany do diagramu zbiorów morfizmów. Korzystamy z tego, że $\kappa$-filtrowanie implikuje zwykłe filtrowanie.
\end{proof}

\begin{twierdzenie}[Prezentowalność zbiorów]\label{cat:thm:accessible-set-presentability}
Zbiór $S$ jest $\kappa$-prezentowalny w $\Set$ wtedy i tylko wtedy, gdy $|S|<\kappa$.
\end{twierdzenie}
\begin{proof}
Jeżeli $|S|<\kappa$, to $\Hom_{\Set}(S,-)$ jest produktem $|S|$ kopii funktora identycznościowego. Twierdzenie~\ref{cat:thm:accessible-filtered-limits}, zastosowane do dyskretnej kategorii indeksującej $S$, daje zachowywanie $\kappa$-filtrowanych kogranic.

Odwrotnie, przykład~\ref{cat:ex:accessible-small-subsets} przedstawia $S$ jako $\kappa$-filtrowaną kogranicę małych podzbiorów. Jeżeli $S$ jest $\kappa$-prezentowalny, to $\id_S$ faktoryzuje się przez pewną inkluzję $A\hookrightarrow S$ z $|A|<\kappa$. Z równości $\iota_A f=\id_S$ wynika, że każdy element $S$ należy do $A$. Stąd $S=A$ i $|S|<\kappa$.
\end{proof}

\begin{stwierdzenie}[Własności domknięcia]\label{cat:prop:accessible-presentable-closure}
Obiekty $\kappa$-prezentowalne są zamknięte na izomorfizmy, retrakty i istniejące kogranice diagramów o mniej niż $\kappa$ morfizmach.
\end{stwierdzenie}
\begin{proof}
Izomorfizm obiektów indukuje izomorfizm ich funktorów reprezentowalnych. Jeżeli $P$ jest retraktem $Q$, tj. istnieją $s:P\to Q$ i $r:Q\to P$ z $rs=\id_P$, to prekompozycja przez $r$ i $s$ przedstawia $\Hom(P,-)$ jako retrakt $\Hom(Q,-)$. Naturalność porównania~\eqref{cat:eq:accessible-hom-colim} przedstawia porównanie dla $P$ jako retrakt porównania dla $Q$ w kategorii strzałek. Retrakt izomorfizmu jest izomorfizmem: dla morfizmów strzałek $f\to g\to f$ o złożeniu identyczności i odwracalnym $g$, złożenie odpowiednich morfizmów składowych z $g^{-1}$ daje obustronną odwrotność $f$.

Niech teraz $P=\operatorname{colim}_{j\in J}P_j$, gdzie $|\operatorname{Mor}J|<\kappa$ i wszystkie $P_j$ są $\kappa$-prezentowalne. Własność uniwersalna kogranicy daje naturalną bijekcję
\[
 \Hom(P,X)\cong\lim_{j\in J^{\mathrm{op}}}\Hom(P_j,X).
\]
Dla $\kappa$-filtrowanego diagramu $D$ stosujemy tę bijekcję, prezentowalność poszczególnych $P_j$ i twierdzenie~\ref{cat:thm:accessible-filtered-limits}:
\[
\begin{aligned}
 \Hom(P,\operatorname{colim}_iD(i))
 &\cong\lim_j\Hom(P_j,\operatorname{colim}_iD(i))\\
 &\cong\lim_j\operatorname{colim}_i\Hom(P_j,D(i))\\
 &\cong\operatorname{colim}_i\lim_j\Hom(P_j,D(i))\\
 &\cong\operatorname{colim}_i\Hom(P,D(i)).
\end{aligned}
\]
Każda z bijekcji jest indukowana przez odpowiednie morfizmy strukturalne. Ich złożenie jest odwrotnością porównania~\eqref{cat:eq:accessible-hom-colim}. Zatem $P$ jest $\kappa$-prezentowalny.
\end{proof}

\subsection{Moduły i prezentacje przez generatory oraz relacje}

Niech $R$ będzie pierścieniem z jedynką, a $\mathcal C$ kategorią lewych $R$-modułów. Dla zbioru $S$ oznaczamy przez $R^{(S)}$ wolny lewy moduł o bazie $S$.

\begin{lemat}[Kogranice filtrowane modułów]\label{cat:lem:accessible-module-filtered}
Funktor zapominania z kategorii lewych $R$-modułów do $\Set$ zachowuje małe filtrowane kogranice.
\end{lemat}
\begin{proof}
Niech $D:I\to R\text{-}\cat{Mod}$ będzie filtrowanym diagramem. W zbiorowej kogranicy jego zbiorów bazowych oznaczmy klasę elementu $x\in D(i)$ przez $[i,x]$. Dla dwóch klas wybieramy wspólny stopień $k$ z morfizmami $u:i\to k$, $v:j\to k$ i definiujemy
\[
 [i,x]+[j,y]=[k,D(u)x+D(v)y],\qquad
 r[i,x]=[i,rx].
\]
Klasa zera nie zależy od stopnia. Niezależność od reprezentantów i wspólnego stopnia wynika z lematu~\ref{cat:lem:filtered-equality} i filtrowania: skończenie wiele użytych morfizmów oraz równości można przenieść do jednego stopnia, gdzie reprezentanty są równe. Tam dodawanie i mnożenie przez skalar są dobrze określone, ponieważ morfizmy diagramu są liniowe. W ten sam sposób każdy aksjomat modułu sprowadza się do aksjomatu w jednym stopniu, gdyż dotyczy skończenie wielu elementów i działań. Morfizmy strukturalne są liniowe przez konstrukcję. Dla zgodnej rodziny morfizmów liniowych $h_i:D(i)\to M$ jedyna funkcja z kogranicy zbiorowej jest dana przez $[i,x]\mapsto h_i(x)$. Jest liniowa na mocy wzorów na działania i zgodności rodziny. Dowodzi to własności uniwersalnej kogranicy modułów.
\end{proof}

\begin{twierdzenie}[Prezentowalność modułów]\label{cat:thm:accessible-module-presentability}
Lewy $R$-moduł $M$ jest $\kappa$-prezentowalny wtedy i tylko wtedy, gdy istnieje ciąg dokładny
\[
 R^{(T)}\xrightarrow{d}R^{(S)}\longrightarrow M\longrightarrow0
 \qquad(|S|,|T|<\kappa).
\]
\end{twierdzenie}
\begin{proof}
Jeżeli istnieje taka prezentacja, to dla dowolnego $N$ zbiór $\Hom_R(M,N)$ jest equalizatorem dwóch funkcji
\[
 \Hom_R(R^{(S)},N)\rightrightarrows\Hom_R(R^{(T)},N),
\]
pierwszej zadanej przez prekompozycję z $d$, drugiej stale równej zeru. Z własności wolnego modułu oba zbiory po prawej są produktami odpowiednio $|S|$ i $|T|$ kopii zbioru bazowego $N$. Lemat~\ref{cat:lem:accessible-module-filtered} oraz twierdzenie~\ref{cat:thm:accessible-filtered-limits} dają zachowywanie $\kappa$-filtrowanych kogranic przez te produkty oraz equalizator. Otrzymujemy prezentowalność $M$.

Odwrotnie, najpierw przedstawmy dowolny moduł jako $\kappa$-filtrowaną kogranicę modułów o małych prezentacjach. Wybieramy surjekcję $F=R^{(S_0)}\to M$, np. z bazą indeksowaną przez elementy $M$, i oznaczamy jej jądro przez $K$. Niech $I$ będzie posetem par $(S,T)$, gdzie $S\subseteq S_0$, $|S|<\kappa$, a $T\subseteq K\cap R^{(S)}$ i $|T|<\kappa$. Porządek jest zadany przez zawieranie w obu współrzędnych. Regularność daje $\kappa$-filtrowanie: wspólnym ograniczeniem górnym rodziny mniej niż $\kappa$ par jest para sum ich współrzędnych. Definiujemy
\[
 M_{S,T}=R^{(S)}/\langle T\rangle_R .
\]
Inkluzje indeksów indukują dobrze określone morfizmy ilorazowe oraz zgodną rodzinę morfizmów do $M$.

Każdy element $M$ ma reprezentanta o skończonym nośniku w $F$, więc pochodzi z pewnego $M_{S,T}$. Jeśli dwa reprezentanty mają ten sam obraz w $M$, to po przeniesieniu do wspólnego stopnia ich różnica należy do $K$. Dołączając tę różnicę do zbioru relacji $T$, otrzymujemy większy stopień, w którym reprezentanty są równe. Dodanie jednej relacji nie przekracza ograniczenia kardynalnego. Lemat~\ref{cat:lem:accessible-module-filtered} dowodzi więc, że kanoniczny morfizm $\operatorname{colim}_I M_{S,T}\to M$ jest bijekcją na zbiorach bazowych, a zatem izomorfizmem modułów.

Jeśli $M$ jest $\kappa$-prezentowalny, faktoryzacja jego identyczności daje morfizmy $s:M\to Q$ oraz $r:Q\to M$ z $rs=\id_M$, gdzie $Q$ jest jednym z modułów $M_{S,T}$. Pozostaje wykazać, że retrakt modułu o małej prezentacji ma małą prezentację. Niech $F_1\to F_0\xrightarrow{\pi}Q\to0$ będzie taką prezentacją. Endomorfizm $e=sr$ modułu $Q$ jest idempotentem. Mamy
\[
 \ker r=\operatorname{im}(\id_Q-e):
\]
jedna inkluzja wynika z $r(\id_Q-e)=0$, a druga z $(\id_Q-e)q=q$ dla $rq=0$. Moduł $Q$ ma mniej niż $\kappa$ generatorów, więc $\ker r$ jest generowane przez obrazy tych generatorów pod $\id_Q-e$. Wybieramy po jednym podniesieniu tych obrazów do $F_0$. Jądro $r\pi$ jest generowane przez te podniesienia razem z generatorami $\ker\pi$: dla elementu $x\in\ker(r\pi)$ odejmujemy odpowiednią kombinację podniesień, aby otrzymać element $\ker\pi$. Obie rodziny generatorów mają moc mniejszą od $\kappa$, więc ich suma również. Surjekcja $r\pi:F_0\to M$ daje żądaną prezentację.
\end{proof}

\begin{przyklad}[Generatory nie wystarczają]\label{cat:ex:accessible-fg-not-fp}
Niech $R=k[x_1,x_2,\ldots]$ dla ciała $k$, a $J=(x_1,x_2,\ldots)$. Moduł $R/J$ jest generowany przez klasę $1$, ale nie jest skończenie prezentowalny. Ideał $J$ nie jest skończenie generowany: dowolne skończenie wiele jego elementów zależy od skończenie wielu zmiennych i zeruje się po podstawieniu tych zmiennych równych zeru, podczas gdy pewna niewystępująca zmienna nie zeruje się przy takim podstawieniu.

Gdyby $R/J$ miał skończoną prezentację $R^m\to R^n\xrightarrow{q}R/J\to0$, wybieramy morfizm $a:R^n\to R$ podnoszący $q$ wzdłuż $\pi:R\to R/J$ oraz $b:R\to R^n$ podnoszący $\pi$ wzdłuż $q$ (wolne moduły są rzutowe). Wtedy $J$ jest generowane przez obrazy pod $a$ skończonego zbioru generatorów $\ker q$ oraz element $(\id_R-ab)(1)$. Istotnie wszystkie te elementy należą do $J$, a dla $x\in J$ mamy $bx\in\ker q$ i
\[
 x=abx+(\id_R-ab)x.
\]
Byłoby to sprzeczne z brakiem skończonego układu generatorów $J$. Twierdzenie~\ref{cat:thm:accessible-module-presentability} pokazuje więc, że $R/J$ nie jest $\aleph_0$-prezentowalny.
\end{przyklad}
